% ============================================================================
% GEM SECTION 2 --- merged method.
%
% EDITS APPLIED HERE ARE ONLY THOSE THE ABRIDGMENT MAP SPECIFIES.  Retained
% prose is VERBATIM from paper_iclr2027/sections/{process,control}.tex.  No
% sentence has been reworded to save space; page fitting is left to the author.
%
%   2.1 <- merge of 4.1-4.3.  Map's MOVE list, applied exactly:
%          the formal union X = |_| X_n; node/edge vocabularies; the detailed
%          valence and connectivity guards; Proposition 1 + proof; the mark
%          factorization; the teacher-program construction; the training
%          objective derivation; the generator-matching divergence discussion;
%          Proposition 2 + proof.
%   2.2 <- merge of 4.4-4.5.  Map keeps ONLY the h_b interpretation, the
%          scalable controlled kernel, and the two-sentence follow-up.  The
%          full recursion, theorem, proof, Monte Carlo target, Bellman
%          auxiliary loss and exact-vs-approximate qualifications MOVE.
%   2.3 <- NEW, one paragraph, supplied verbatim by the map.
%   MOVED WHOLE: 4.6 SMC, 4.7 preference-conditioned control, 4.8 retargeting
%          and pathwise constraints.
% ============================================================================
\section{Executable Molecular Rewriting and Control}
\label{sec:method}

We develop \method{} as a reusable stochastic process over molecular edits and a finite-horizon controller over that process. \method{} realizes generative stochastic rewriting by repeatedly applying executable molecular edits under a learned, goal-independent transition law. An exact rewrite system first determines which transitions between complete molecular graphs are executable. We then learn a goal-independent reference law over this state-dependent support and push it forward to a canonical molecular successor kernel $R_\theta$. Finally, we control the frozen successor process through a budget-conditioned $h$-transform, using a learned future-value model $h_\phi$ to approximate the required backward values at molecular scale.

\subsection{Executable reference process}
\label{sec:gem-support}

\method{} represents every generative state as a complete molecular graph, drawn from a state space $\X$ decomposed by heavy-atom count (\cref{app:process}). Thus each state $x\in\X$ is a complete molecule rather than a masked, corrupted, or partially specified representation.

Molecular transitions are generated by typed graph rewrites. An edit mark $a$ specifies a rewrite type together with its molecular match and typed payload, and $\T(x,a)$ denotes the graph obtained by applying that rewrite to $x$. The full atom-, bond-, and ring-level edit language is specified in \cref{app:operators}. From a state $x$, the executor exposes only the finite, state-dependent set
\begin{equation}
\A(x)=\bigl\{\,a:\ a\ \text{matches}\ x\ \ \text{and}\ \ \T(x,a)\ \text{passes all executor guards}\,\bigr\}.
\label{eq:gem-editset}
\end{equation}
Execution follows a propose--validate--commit rule: a successor is admitted only if its attributes remain in the declared vocabularies, affected valences are satisfied, the graph remains connected, and its heavy-atom count lies in $\{1,\dots,n_{\max}\}$. An edit that fails any guard contributes no transition to $\A(x)$. Inadmissible molecular states are therefore absent from the stochastic support rather than discouraged through a learned penalty.

The decomposition of $\X$ by heavy-atom count makes the process trans-dimensional. Each admitted edit changes heavy-atom count by at most one, yielding \emph{birth}, \emph{death}, or \emph{same-cardinality} transitions between $\X_n$ and $\X_{n+1}$, $\X_{n-1}$, or $\X_n$, respectively. A trajectory may therefore grow, shrink, and restructure a molecular graph within a common state space. Every state reached by a process supported on $\{\A(x)\}_{x\in\X}$ remains in the declared molecular state space (\cref{app:proofs}).

Multiple legal marks can produce the same molecular graph because of graph symmetries, equivalent matches, or alternative internal rewrite descriptions. We therefore push the learned mark law through the executor and aggregate probability over all marks that produce the same canonical molecular graph:
\begin{equation}
R_\theta(y\mid x)\ \propto\!\!\sum_{a:\,\T(x,a)\simeq y}\!\! p_\theta(a\mid x),
\label{eq:gem-rtheta}
\end{equation}
where $\simeq$ denotes equivalence under the registered molecular canonicalization and the normalization runs over the distinct non-self canonical successors $\Sset(x)$. $R_\theta$ is trained by canonical-successor negative log-likelihood on executable teacher transitions and receives no design objective or downstream task information. The executor therefore defines possibility, while $R_\theta$ learns a state-dependent probability over molecular outcomes. Control acts after canonicalization; reweighting individual marks beforehand can reintroduce dependence on arbitrary edit multiplicities (\cref{app:quotient}).

\subsection{Future-aware molecular control}
\label{sec:gem-control}

We introduce task-specific control without modifying the reference process. Let $z$ denote a design objective, $g_z:\X\rightarrow\R_{\ge0}$ its terminal desirability, and $h_b(x,z)$ the backward value at remaining budget $b$. Thus $h_b(x,z)$ evaluates a molecule by the desirability reachable from it after $b$ further transitions under the frozen reference process. Unlike a score of the current molecule, it depends on the molecular futures that remain accessible through $R_\theta$.

Replacing the exact backward value with a learned goal- and budget-conditioned approximation $h_\phi$ yields the scalable controlled kernel
\begin{equation}
P_\phi(y\mid x,z,b)=\frac{R_\theta(y\mid x)\,h_\phi(y,z,b-1)}{\displaystyle\sum_{y'\in\Sset(x)}R_\theta(y'\mid x)\,h_\phi(y',z,b-1)},
\qquad y\in\Sset(x).
\label{eq:gem-central}
\end{equation}
The controller therefore reweights only executable successors already represented by $R_\theta$. With exact backward values this recovers the finite-horizon Doob transform; at molecular scale, $h_\phi$ is learned from finite-budget continuations sampled from the frozen reference process. The backward recursion, the exact-control theorem and its proof, the Monte Carlo and Bellman training targets, the exact-versus-approximate qualifications, and the sequential Monte Carlo sampler are given in \cref{app:control,app:proofs,app:controller}.

\subsection{Temporally extended executable transformations}
\label{sec:gem-macros}

Useful molecular transformations may require several primitive edits before their value is realized. \method{} therefore permits temporally extended programs, including linked and fused ring-system construction, that compile entirely into the same executable atom- and bond-level rewrites. These programs organize and restrict existing molecular support rather than introducing fragment templates or new chemistry: the controller selects transformation semantics, while $R_\theta$ ranks legal realizations. For endpoint-constrained tasks with an expensive oracle, initially infeasible seeds can first use an oracle-free, diversity-preserving frontier over the same executable state space; expensive objective evaluation begins only after feasible molecular states are reached. Implementation, macro qualification, ring composition, frontier construction, and lineage tracking are given in \cref{app:macros}.
